\newpage
\section[PHT3D Exercise 5: Pyrite oxidation during deep well injection]{PHT3D Exercise 5: Pyrite oxidation during deep well injection}

Managed aquifer recharge is increasingly used to enhance the sustainable development of water supplies. Common recharge techniques include aquifer storage and recovery (ASR), infiltration ponds, river bank filtration and deep-well injection. Following recharge the water quality of the injectant is typically altered by a multitude of geochemical processes during subsurface passage and storage. Relevant geochemical processes that affect the major ion chemistry include microbially mediated redox reactions, mineral dissolution/precipitation, sorption and ion-exchange. The hydrochemical conditions and changes that occur under these circumstances, in particular the temporal and spatial changes of pH and redox conditions, are in many cases the controlling factor for the fate of micropollutants such as herbicides and pharmaceuticals. Similarly, changes in mineralogical composition such as dissolution and precipitation of iron- or aluminiumoxides may affect the mobility of trace metals as well as the attachment and subsequent decay of pathogenic viruses. Laboratory and field-scale experimental studies are aimed at investigating such processes under controlled conditions and to eventually develop a better qualitative and quantitative understanding of their complex interactions, both site-specific and at a fundamental level. \
\cite{prommer_stuyfzand_2005} carried out a reactive transport 
modelling study to analyse the data collected during a deep well 
injection experiment in an anaerobic, pyritic aquifer near Someren in Southern Netherlands. 
